\documentclass[12pt,a4paper]{article}

\usepackage[utf8]{inputenc}
\usepackage[T2A]{fontenc}
\usepackage[russian]{babel}
\usepackage{amsmath,amssymb}
\usepackage{geometry}
\usepackage{graphicx}
\usepackage{float}

\geometry{margin=2.2cm}

\newcommand{\plotplaceholder}[2][0.9\textwidth]{%
    \IfFileExists{#2}{%
        \includegraphics[width=#1]{#2}%
    }{%
        \fbox{%
            \begin{minipage}[c][0.22\textheight][c]{#1}
            \centering
            Заглушка для графика\\[4pt]
            \texttt{\detokenize{#2}}
            \end{minipage}%
        }%
    }%
}

\begin{document}

\title{Графики}
\author{}
\date{}
\maketitle

\section*{Графическая иллюстрация аналитически полученных результатов}

Основные графики построены для $n=30$ и $n=100$ по формуле
$\bar{\delta}_n(\xi)$ без дополнительных аппроксимаций.

\begin{figure}[H]
\centering
\plotplaceholder[0.82\textwidth]{delta_bar_n30_100.png}
\caption{Графики ОИСКО $\bar{\delta}_{30}(\xi)$ и $\bar{\delta}_{100}(\xi)$, построенные по аналитической формуле для прямоугольного ядра.}
\end{figure}

При малых $\xi$ ошибка быстро растёт из-за дисперсионного члена
$\sim (2an\xi)^{-1}$, а при больших $\xi$ кривая приближается к пределу
$1/(2\sqrt{3})$.

\begin{figure}[H]
\centering
\plotplaceholder[0.82\textwidth]{efficiency_en_n30_100.png}
\caption{Эффективности $e_{30}(\xi)$ и $e_{100}(\xi)$ с максимумом, равным $1$.}
\end{figure}

Эффективность максимальна вблизи оптимального значения $\xi$ и убывает как
при слишком малых, так и при слишком больших значениях параметра.

\begin{figure}[H]
\centering
\plotplaceholder[0.82\textwidth]{xi_opt_vs_n.png}
\caption{Дополнительно: зависимость $\xi_{\mathrm{opt}}(n)$ от размера выборки $n$.}
\end{figure}

По мере роста объёма выборки оптимальное значение $\xi$ убывает, поскольку
для больших $n$ можно использовать более узкое окно без сильного штрафа по
дисперсии.

\begin{figure}[H]
\centering
\plotplaceholder[0.82\textwidth]{delta_min_vs_n.png}
\caption{Дополнительно: зависимость $\bar{\delta}_{n,\min}$ от размера выборки $n$.}
\end{figure}

Минимальная ОИСКО монотонно убывает: увеличение выборки улучшает наилучшее
достижимое качество оценки.

\end{document}
